% Part: many-valued-logic
% Chapter: syntax-and-semantics
% Section: connectives

\documentclass[../../../include/open-logic-section]{subfiles}

\begin{document}

\olfileid{mvl}{syn}{con}

\olsection{زبان‌ها و رابط‌ها}

منطق گزاره‌ای کلاسیک و بسیاری منطق‌های دیگر از مجموعه‌ای ثابت از
\emph{ثابت‌های گزاره‌ای} و \emph{رابط‌ها} استفاده می‌کنند. برای نمونه،
موارد زیر را نمادهای اولیه می‌گیریم:
\begin{enumerate}
\tagitem{prvFalse}{ثابت گزاره‌ای~$\lfalse$ برای~!!{falsity}.}{}
\tagitem{prvTrue}{ثابت گزاره‌ای~$\ltrue$ برای~!!{truth}.}{}
\item رابط‌های منطقی:
  \startycommalist
  \iftag{prvNot}{\ycomma $\lnot$ (نقیض)}{}%
  \iftag{prvAnd}{\ycomma $\land$ (عطف)}{}%
  \iftag{prvOr}{\ycomma $\lor$ (فصل)}{}%
  \iftag{prvIf}{\ycomma $\lif$ (!!{conditional})}{}%
  \iftag{prvIff}{\ycomma $\liff$ (!!{biconditional})}{}%
\end{enumerate}
\iftag{defNot,defOr,defAnd,defIf,defIff,defTrue,defFalse,defEx,defAll}{%
افزون بر رابط‌های اولیهٔ بالا، از نمادهایی نیز استفاده می‌کنیم که
به‌صورت اختصار تعریف شده‌اند، مانند
\startycommalist
  \iftag{defNot}{\ycomma $\lnot$ (نقیض)}{}%
  \iftag{defAnd}{\ycomma $\land$ (عطف)}{}%
  \iftag{defOr}{\ycomma $\lor$ (فصل)}{}%
  \iftag{defIf}{\ycomma $\lif$ (!!{conditional})}{}%
  \iftag{defIff}{\ycomma $\liff$ (!!{biconditional})}{}%
  \iftag{defFalse}{\ycomma $\lfalse$ (!!{falsity})}{}%
  \iftag{defTrue}{\ycomma $\ltrue$ (!!{truth})}.}{}

همین رابط‌ها در منطق‌های چندارزشی نیز به کار می‌روند. بااین‌حال، اغلب
سودمند است که نسخه‌های متفاوتی مثلاً از عطف را در یک منطق بگنجانیم؛
در آن صورت برای متمایز نگه‌داشتن نسخه‌ها به نمادهای متفاوت نیاز داریم.
برخی منطق‌های چندارزشی رابط‌هایی نیز دارند که در منطق کلاسیک هیچ همتایی
ندارند. بنابراین کمی کلی‌تر از معمول عمل می‌کنیم.

\begin{defn}
  یک \emph{زبان گزاره‌ای} از مجموعهٔ~$\Lang L$ از \emph{رابط‌ها} تشکیل
  می‌شود. هر رابط~$\star$ یک \emph{تعداد موضع} دارد؛ رابطی با تعداد
  موضع~$n$ را \emph{$n$موضعی} می‌نامیم. رابط‌های صفرموضعی را
  \emph{ثابت}، رابط‌های یک‌موضعی را \emph{یکانی} و رابط‌های دوموضعی را
  \emph{دوتایی} نیز می‌نامیم.
\end{defn}

\begin{ex}
  زبان استاندارد منطق گزاره‌ای~$\Lang L_0$ از رابط‌های زیر، با تعداد
  موضع متناظر، تشکیل می‌شود:
  $\lfalse$~($0$)،
  $\lnot$~($1$)،
  $\land$~($2$)،
  $\lor$~($2$)،
  $\lif$~($2$). بیشتر منطق‌هایی که بررسی می‌کنیم از این زبان استفاده
  خواهند کرد. بعضی منطق‌ها بنا بر سنت و قرارداد برای برخی رابط‌ها
  نمادهای دیگری به کار می‌برند. برای نمونه، در منطق ضربی معمولاً به‌جای
  $\land$ از نماد عطف~$\odot$ استفاده می‌شود. گاهی نیز افزودن عملگر
  تازه‌ای مانند عملگر تعین~$\triangle$ (یک‌موضعی) سودمند است.
\end{ex}

\end{document}
